\section{Conclusion}
\label{sec:conclusion}

Optimizer selection for large-scale model training has become a system-level,
multi-objective decision, yet the field offers more than one hundred methods
described in incompatible vocabularies and supported by protocol-sensitive
evidence. This paper set out to make this body of work navigable and, ultimately,
actionable. We treated every optimizer as a structured transformation of a
stochastic training signal and showed that modern methods are sparse
modifications of a shared five-stage meta-pipeline. We then used norm-constrained linear minimization oracles to place sign,
spectral-orthogonalized,
Kronecker-factored preconditioned, and projected updates inside a single four-axis
decomposition. Finally, we organized the resulting mechanisms into a
dual-dimension taxonomy that records both where a method acts and which training
objective it targets. A controlled benchmark over language-model pretraining
then turned this framework from a descriptive device into an evaluative one. The benchmark spans four scales from 60M to 1B parameters, four architectures,
and context lengths from 256 to 32k tokens, with a targeted CIFAR100
vision-backbone study used to probe architecture-dependent optimizer behavior
beyond language modeling.

Our main conclusion is simple. There is no universal best optimizer, and the
practical question is which mechanism's strength matches the binding constraint
of the training regime. Under the
tested protocol, the strongest quality gains are associated with
geometry-sensitive direction maps and the structured state that feeds them.
Scalar refinements of the adaptive-moment template alone contribute little. Most
element-wise variants of AdamW do not survive a retuned baseline. Two empirical
regularities sharpen this picture and should be treated as design constraints.
Aggressive state compression is rank-bounded and
degrades sharply as the gradient's effective rank rises with context length.
Spectral matrix geometry is architecture-conditional, so its gains must be
validated on the target topology, and their transfer must be confirmed.
Composability, in turn, is governed by locality on the four axes. Techniques on
different axes or pipeline stages stack, while techniques that contend for the
same slot require an explicit ordering. The unit of design is therefore a
composition chosen against a budget.

These findings translate into concrete guidance. Across the benchmarked
methods, the mainstream choices and the regime each one fits
(Table~\ref{tab:optimizer_tiers}) are the following.
\begin{itemize}
  \item \textbf{AdamW} is the default reference for general-purpose
  pretraining. It is not the strongest on any single metric, but it is stable,
  inexpensive, interpretable, and the baseline every other choice should be
  measured against.
  \item \textbf{RMNP}~\citep{deng2026rmnp} is better read as a
  practical, architecture-aware step beyond standard Muon-style preconditioning.
  It targets a better quality--efficiency balance by
  using architecture-induced row-wise structure. When the relevant curvature and
  momentum statistics are diagonally dominant in a row-wise sense, a Muon-style
  matrix direction can be implemented through a cheaper row-normalized form.
  \item \textbf{SOAP} is the quality ceiling for long-context training and the
  most transferable method across architectures. It is too costly to be a
  default, but it is useful when final quality dominates and compute and memory
  are not the bottleneck.
  \item \textbf{Muon} is a stable, mechanistically transparent
  matrix method. It should be used with awareness of the target model topology.
  \item \textbf{AdaFactor} is the safe low-memory baseline at moderate quality.
  \textbf{APOLLO} is a high-reward, high-risk compressed-state method that excels
  at short context but collapses at long context, so its short-context wins do
  not certify it.
  \item \textbf{Lion} is a cheap exploratory option with an expected quality gap.
  Curvature-aware and geometry-regularized methods such as Sophia and LAMB remain
  situational under the current protocol.
\end{itemize}
In short, the decision rule is to identify the single binding constraint of a
run, such as stability, quality, runtime, memory, or cross-scenario transfer,
and select the method whose dominant strength addresses it. Starting from AdamW,
move to RMNP, SOAP, or a memory-efficient method only when a specific constraint
demands it.

Beyond these recommendations, the framework's main contribution to the
community is an operational coordinate system. Any optimizer, including methods
that have not yet been proposed, can be located by its pipeline stages and its
four-axis coordinate. Its composability with other techniques can be predicted from that
location, and competing methods can be compared under
explicit mechanism and objective assumptions.
This gives practitioners a shared vocabulary and a constraint-driven decision
procedure. For method designers, it maps crowded regions, unexplored
compositions, and empirical boundaries that any new method must respect, such as
rank-bounded compression and architecture-conditional geometry. A further
lesson, developed in Section~\ref{sec:discussion-open-problems}, is that
optimizer geometry should be coupled to the model architecture and training
dynamics, with RMNP as one concrete example. The
same coordinate system marks where this study should grow next. One direction is to develop
quantitative interpretability metrics that turn the present qualitative mechanism
attributions into predictive diagnostics. A
second direction is to extend the study beyond language modeling, building on the
foundation established here.
